\documentclass[10pt,a4paper]{amsart}
\usepackage[margin=19mm]{geometry}
\usepackage{amsmath,amssymb,mathtools}
\usepackage[scr=rsfs,cal=euler]{mathalfa}
\usepackage{xfrac}
\usepackage[unicode,hidelinks,bookmarks=false]{hyperref}
\newcommand{\ie}{i.e.\ }
\newcommand{\cf}{cf.\ }
\newcommand{\resp}{respectively\ }
\newcommand{\Subsection}{\subsection}
\providecommand{\nobreakdash}{\nobreak\hbox{-}}
\setlength{\emergencystretch}{2em}
\multlinegap0pt
\usepackage{xcolor}
\usepackage{amssymb}
\newtheorem{theorem}{Theorem}
\title{Code distances: a new family of invariants of linear codes}
\author{Eduardo Camps-Moreno, Elisa Gorla, Hiram H. L\'opez}
\date{Source: arXiv:2509.16424v2 (CC BY 4.0)}
\begin{document}
\maketitle

\noindent\emph{Source and licence.} Code distances: a new family of invariants of linear codes. Version arXiv:2509.16424v2 (CC BY 4.0), distributed by arXiv under the Creative Commons Attribution 4.0 International licence, http://creativecommons.org/licenses/by/4.0/, as recorded in the arXiv licence metadata for this exact version and shown on the abstract page https://arxiv.org/abs/2509.16424v2. Full licence notice: \texttt{licenses/arxiv-2509.16424-CC-BY-4.0.txt}.\par\smallskip

\noindent\textit{Editorial scope.} Counted unit: the article's theorem theorem (source0.tex lines 818--821). The article's complete original proof of it is delivered with it (source0.tex lines 823--844, one \texttt{proof} environment of 22 lines). No mathematics is written, repaired or re-derived: the statement and the proof are the article's own lines, in the article's own notation, and no retained line is substituted or re-flowed.  No further source-local context is needed: every cross-reference of every retained line resolves inside the two delivered spans.  Nothing was omitted from any retained span, so the package carries no omission record.  No retained line contains a citation, so the wrapper carries no bibliography and no bibliography entry.  No retained line contains a figure environment, an image-inclusion command or drawing code, so no image is delivered and the package declares no asset dependency.  Editorial formatting only, in addition: an AMS \texttt{amsart} preamble that restates the article's own declarations and macros used by the retained lines, namely the environments \texttt{theorem} on separate counters, together with the standard packages those lines need.  The retained environments print in this document as Theorem 1 in order of appearance, whereas the article prints the same items with the numbering of its own full document.  The complete native source of the article as distributed in the arXiv e-print for this exact version is bundled unchanged as \texttt{sources/185376/source0.tex}.
\begin{theorem}
 Let $C\subseteq\mathbb{F}_q^n$. For any $1\leq i\leq k$, there is $t\geq 1$ such that for any $\ell\geq t$, 
 $$\alpha_i^\ell(C)=\alpha_i^\infty(C).$$
\end{theorem}

\begin{proof}
Let $1\leq i\leq k$ and $j=k-i$ and assume $\alpha_i^\infty(C)=d$. Let $H$ be a parity check matrix of $C$ and let $$\mathbf{x}=\begin{pmatrix} x_{11}&\cdots&x_{1n}\\ \vdots&\ddots&\vdots\\
    x_{j1}&\cdots&x_{jn}\end{pmatrix}.$$

Let $T(\mathbf{x})=\begin{pmatrix} H\\ \mathbf{x}\end{pmatrix}$ and define the following sets of ideals in $\bar{\mathbb{F}_q}[\mathbf{x}]$:
\begin{itemize}
\item $I_{dim}$ the ideal of minors of size $n-i$ of $T(\mathbf{x})$.
\item For each $J\subseteq [n]$ of size $d-1$, let $I_J$ be the ideal of minors of size $d-1$ of the submatrix of $T(\mathbf{x})$ whose columns are indexed by $J$.
\item $$I_{C,i}=I_{\dim}\cap\bigcap_{J\subseteq[n], |J|=d-1} I_J.$$
\end{itemize}

Let $P\in\mathbb{F}_q^{jn}$. The previous ideals provide us with the following information:
\begin{itemize}
\item If $P\notin Z(I_{\dim})$, then $T(P)$ has rank $n-i$ and it defines a parity check for a subcode of $C$ of dimension $i$.
\item If for any $J\subseteq [n]$ of size $d-1$, $P\notin Z(I_J)$, then any $d-1$ columns of $T(P)$ are linearly independent, hence the code with parity check matrix $T(P)$ has minimum distance at least $d$.
\item By the previous two items, if $P\notin Z(I_{C,i})$, then $T(P)$ is the parity check matrix of a subcode of $C$ of dimension $i$ and minimum distance greater than or equal to $d$. Equality follows from the maximality of $\alpha_i^\infty(C)$.
\end{itemize}

Let $\mu(I_{C,i})$ be the minimum degree of an element of $I_{C,i}$ and let $$t=\min\{r\ :\ \mu(I_{C,i})< q^r\}.$$

Then for any $\ell\geq t$ one has $I_{C,i}\subseteq\sqrt {I_{C,i}}\not\subseteq (x_{gh}^{q^\ell}-x_{gh})_{g,h=1}^{j,n}$, hence $\mathbb{F}_{q^\ell}^n\setminus Z(I_{C,i})\neq \emptyset$. This guarantees that for any $\ell\geq t$, $\alpha_i^\ell(C)=d$.
\end{proof}

\end{document}
